\documentclass[aspectratio=169,handout,11pt]{beamer}

% =========================
% PAQUETES
% =========================
\usepackage[utf8]{inputenc}
\usepackage[T1]{fontenc}
\usepackage[spanish,es-noshorthands]{babel}
\usepackage{amsmath,amssymb,mathtools}
\usepackage{lmodern}
\usepackage{xcolor}

% =========================
% TEMA
% =========================
\usetheme{Madrid}
\setbeamertemplate{navigation symbols}{}

% =========================
% COLORES
% =========================
\definecolor{AgroVerde}{RGB}{34,102,51}
\definecolor{AgroCrema}{RGB}{248,245,238}
\definecolor{AgroAzul}{RGB}{20,90,130}
\definecolor{AgroRojo}{RGB}{170,50,40}

\setbeamercolor{title}{fg=white,bg=AgroVerde}
\setbeamercolor{frametitle}{fg=white,bg=AgroVerde}
\setbeamercolor{block title}{fg=white,bg=AgroVerde}
\setbeamercolor{block body}{fg=black,bg=AgroCrema}
\setbeamercolor{block title alerted}{fg=white,bg=AgroRojo}
\setbeamercolor{block body alerted}{fg=black,bg=red!5}
\setbeamercolor{block title example}{fg=white,bg=AgroAzul}
\setbeamercolor{block body example}{fg=black,bg=blue!5}

% =========================
% ENTORNOS
% =========================
\newenvironment{idea}{\begin{block}{Idea clave}}{\end{block}}
\newenvironment{ejemplo}{\begin{exampleblock}{Caso agronómico}}{\end{exampleblock}}
\newenvironment{alerta}{\begin{alertblock}{Alerta técnica}}{\end{alertblock}}

% =========================
% DATOS
% =========================
\title{Fracciones Algebraicas en Agronomía}
\author{Luis Tulio González}
\date{2026}

\begin{document}
	
	% =========================
	\begin{frame}
		\titlepage
	\end{frame}
	
	% =========================
	\begin{frame}{Motivación}
		\begin{idea}
			Las fracciones algebraicas modelan relaciones entre variables agrícolas, pero siempre deben respetar restricciones.
		\end{idea}
		
		\pause
		
		\begin{itemize}[<+->]
			\item Dosis (kg/ha)
			\item Riego (L/h)
			\item Rendimiento (ton/ha)
		\end{itemize}
	\end{frame}
	
	% =========================
	\begin{frame}{Simplificación con restricciones}
		\begin{ejemplo}
			\[
			\frac{x^2-9}{x+3}
			\]
		\end{ejemplo}
		
		\pause
		
		\[
		= x-3
		\]
		
		\pause
		
		\begin{block}{Restricción}
			\[
			x \neq -3
			\]
		\end{block}
		
		\pause
		
		\begin{alerta}
			Aunque se simplifique, la restricción se mantiene.
		\end{alerta}
	\end{frame}
	
	% =========================
	\begin{frame}{Ejemplo aplicado}
		\begin{ejemplo}
			\[
			\frac{x^2+5x+6}{x+2}
			\]
		\end{ejemplo}
		
		\pause
		
		\[
		= x+3
		\]
		
		\pause
		
		\begin{block}{Restricción}
			\[
			x \neq -2
			\]
		\end{block}
	\end{frame}
	
	% =========================
	\begin{frame}{Suma con restricciones}
		\begin{ejemplo}
			\[
			\frac{2}{x}+\frac{3}{x+1}
			\]
		\end{ejemplo}
		
		\pause
		
		\[
		= \frac{5x+2}{x(x+1)}
		\]
		
		\pause
		
		\begin{block}{Restricciones}
			\[
			x \neq 0, \quad x \neq -1
			\]
		\end{block}
	\end{frame}
	
	% =========================
	\begin{frame}{Multiplicación}
		\begin{ejemplo}
			\[
			\frac{x^2-9}{x^2-4x+4}\cdot\frac{x-2}{x+3}
			\]
		\end{ejemplo}
		
		\pause
		
		\[
		= \frac{x-3}{x-2}
		\]
		
		\pause
		
		\begin{block}{Restricciones}
			\[
			x \neq 2, \quad x \neq -3
			\]
		\end{block}
	\end{frame}
	
	% =========================
	\begin{frame}{División}
		\begin{ejemplo}
			\[
			\frac{x^2-1}{x}\div\frac{x+1}{x^2}
			\]
		\end{ejemplo}
		
		\pause
		
		\[
		= (x-1)x
		\]
		
		\pause
		
		\begin{block}{Restricciones}
			\[
			x \neq 0, \quad x \neq -1
			\]
		\end{block}
	\end{frame}
	
	% =========================
	\begin{frame}{Errores frecuentes}
		\begin{alerta}
			Errores comunes:
		\end{alerta}
		
		\begin{itemize}[<+->]
			\item Olvidar restricciones
			\item Simplificar mal
			\item Aplicar valores prohibidos
		\end{itemize}
	\end{frame}
	
	% =========================
	\begin{frame}{Ejercicios}
		\small
		
		\begin{enumerate}
			\item \[
			\frac{x^2-25}{x-5}
			\]
			Restricción: \(x \neq 5\)
			
			\item \[
			\frac{2}{x}+\frac{5}{x-3}
			\]
			Restricciones: \(x \neq 0, 3\)
			
			\item \[
			\frac{x^2-9}{x^2-6x+9}
			\]
			Restricción: \(x \neq 3\)
			
		\end{enumerate}
	\end{frame}
	
	% =========================
	\begin{frame}{Ejercicios II}
		\small
		
		\begin{enumerate}
			\setcounter{enumi}{3}
			
			\item \[
			\frac{x+1}{x-2}-\frac{2x}{x+3}
			\]
			Restricciones: \(x \neq 2, -3\)
			
			\item \[
			\frac{x^2-4}{x}\div\frac{x-2}{x+1}
			\]
			Restricciones: \(x \neq 0, 2, -1\)
			
			\item \[
			\frac{3}{x+1}+\frac{4}{x-2}
			\]
			Restricciones: \(x \neq -1, 2\)
			
		\end{enumerate}
	\end{frame}
	
	% =========================
	\begin{frame}{Problema Agronómico}
		\small
		
		\[
		\frac{x^2+7x+12}{x+3}
		\]
		
		\begin{itemize}
			\item Simplifique
			\item Restricción: \(x \neq -3\)
			\item Interprete el resultado
		\end{itemize}
		
	\end{frame}
	
	% =========================
	\begin{frame}{Cierre}
		\begin{idea}
			Toda expresión algebraica tiene restricciones que deben respetarse en la realidad.
		\end{idea}
		
		\pause
		
		\centering
		\textbf{Nunca evaluar donde el modelo no existe}
		
	\end{frame}
	
\end{document}